\chapter{Implementation}

\section{Technology Stack}

The artefact is implemented using the following technologies:

\begin{itemize}
    \item \textbf{Infrastructure as Code:} AWS SAM (Serverless Application Model) with CloudFormation
    \item \textbf{Runtime:} Python 3.12 for all Lambda functions
    \item \textbf{Testing:} pytest 9.1, moto 5.0 (mocked AWS services), ruff (linting)
    \item \textbf{Data Analysis:} pandas 2.2, numpy 2.0, scipy 1.14, scikit-learn 1.5 (Isolation Forest)
    \item \textbf{Visualisation:} matplotlib 3.9
    \item \textbf{AWS Services:} Lambda, API Gateway (REST), DynamoDB, CloudWatch, X-Ray, SSM Parameter Store
    \item \textbf{Baseline Comparison:} RCAEval harness (BARO, CIRCA, TraceRCA, CausalRCA from GitHub)
\end{itemize}

\section{Repository Structure}

\begin{verbatim}
serverless-fault-localisation/
├── template.yaml                 # SAM infrastructure template
├── src/                         # Lambda handlers + shared layer
│   ├── orders_api/app.py
│   ├── inventory/app.py
│   ├── payments/app.py
│   ├── notifications/app.py
│   └── layer/faultlab/          # Shared library
│       ├── fault.py             # Fault switch logic
│       ├── calls.py             # Downstream invocation
│       └── obs.py               # X-Ray patching
├── injector/                    # Fault injection scheduler
│   ├── schedule.py              # Generate injection timeline
│   └── injector.py              # Write faults to SSM
├── detector/                    # Detection & localisation
│   ├── rules.py                 # CloudWatch threshold detector
│   ├── ranker.py                # X-Ray dependency ranking
│   └── collectors.py            # Fetch CloudWatch/X-Ray APIs
├── eval/                        # Evaluation pipeline
│   ├── metrics.py               # P/R/F1, Top-k, delay
│   ├── rig.py                   # End-to-end rig evaluation
│   ├── leg2.py                  # RCAEval comparison
│   ├── overhead.py              # Telemetry volume, latency, cost
│   └── stats.py                 # Hypothesis tests, effect sizes
├── baseline_runner/             # RCAEval harness wrapper
│   ├── rule_arm.py              # Rule-based on RCAEval schema
│   ├── hybrid_arm.py            # IF + trace ranking
│   └── run_baselines.py         # BARO, CIRCA, TraceRCA, CausalRCA
├── workloads/loadgen.py         # Open-loop load generator
├── sim/telemetry.py             # Simulated CloudWatch/X-Ray data
├── configs/
│   ├── experiment.yaml          # Frozen thresholds + hash
│   └── prices.yaml              # AWS list prices (dated)
├── scripts/
│   ├── campaign.py              # Run calibration/steady/peak
│   ├── collect_telemetry.py    # Fetch CloudWatch/X-Ray APIs
│   ├── budget_estimate.py       # Predict campaign costs
│   └── check_bib.py             # Verify DOI resolution
├── docs/                        # Methodology documentation
│   ├── ANALYSIS_PLAN.md
│   ├── CONFIGURATION_MANUAL.md
│   ├── FAULT_TAXONOMY.md
│   └── ETHICS.md
├── tests/                       # 76 unit/integration tests
├── results/                     # Evaluation outputs
│   ├── sim/                     # Simulated rig results
│   └── rcaeval/                 # RCAEval baseline comparison
├── figures/                     # Generated plots
├── bib/references.bib           # Verified bibliography
└── Makefile                     # Build, test, deploy targets
\end{verbatim}

\section{AWS SAM Template Implementation}

The \texttt{template.yaml} defines all infrastructure declaratively:

\subsection{API Gateway REST API}

\begin{verbatim}
OrdersApi:
  Type: AWS::Serverless::Api
  Properties:
    StageName: prod
    TracingEnabled: true   # Enable X-Ray tracing
    MethodSettings:
      - ResourcePath: /*
        HttpMethod: '*'
        LoggingLevel: ERROR
        MetricsEnabled: true
\end{verbatim}

\textbf{Note:} REST API (not HTTP API) is required for X-Ray tracing support. HTTP APIs cannot be traced at the API Gateway level.

\subsection{Lambda Functions}

Each Lambda function is defined with:

\begin{verbatim}
OrdersApiFunction:
  Type: AWS::Serverless::Function
  Properties:
    Runtime: python3.12
    Handler: app.handler
    CodeUri: src/orders_api/
    Timeout: 10
    MemorySize: 256
    Tracing: Active              # Enable X-Ray
    Policies:
      - DynamoDBCrudPolicy:
          TableName: !Ref OrdersTable
      - SSMParameterReadPolicy:
          ParameterName: !Ref FaultParam
    Layers:
      - !Ref FaultLabLayer       # Shared faultlab library
    Environment:
      Variables:
        TABLE_NAME: !Ref OrdersTable
        FAULT_PARAM: !Sub /${AWS::StackName}/fault
        INVENTORY_FN: !Ref InventoryFunction
        PAYMENTS_FN: !Ref PaymentsFunction
        NOTIFICATIONS_FN: !Ref NotificationsFunction
\end{verbatim}

\textbf{Shared Layer:} The \texttt{faultlab} library is packaged as a Lambda layer to avoid code duplication. All functions import \texttt{from faultlab import fault, calls, obs}.

\subsection{DynamoDB Table}

\begin{verbatim}
OrdersTable:
  Type: AWS::DynamoDB::Table
  Properties:
    BillingMode: PAY_PER_REQUEST  # On-demand, no provisioning
    AttributeDefinitions:
      - AttributeName: pk
        AttributeType: S
    KeySchema:
      - AttributeName: pk
        KeyType: HASH
\end{verbatim}

Orders are stored with \texttt{pk = "ORDER\#\{uuid\}"}, inventory SKUs with \texttt{pk = "SKU\#\{sku\}"}.

\subsection{X-Ray Sampling Rule}

\begin{verbatim}
SamplingRule:
  Type: AWS::XRay::SamplingRule
  Properties:
    RuleName: !Sub ${AWS::StackName}-rule
    Priority: 100
    FixedRate: 0.05         # 5% sampling (policy mode)
    ReservoirSize: 0
    ServiceName: '*'
    ServiceType: '*'
    HTTPMethod: '*'
    URLPath: '*'
    ResourceARN: '*'
\end{verbatim}

\textbf{Sampling Control:} For overhead comparison, we deploy twice:
\begin{itemize}
    \item \texttt{FixedRate: 1.0} (full tracing, 100\%)
    \item \texttt{FixedRate: 0.05} (policy mode, 5\%)
\end{itemize}

\subsection{CloudFormation Outputs}

\begin{verbatim}
Outputs:
  ApiUrl:
    Value: !Sub https://${OrdersApi}.execute-api.${AWS::Region}.amazonaws.com/prod
  StackName:
    Value: !Ref AWS::StackName
\end{verbatim}

These outputs are used by \texttt{scripts/campaign.py} to target the deployed stack.

\section{Lambda Handler Implementation}

\subsection{orders-api Handler (Entry Service)}

\begin{verbatim}
def handler(event, context):
    ddb, ssm = clients()
    t0 = time.perf_counter()
    route = f"{event['httpMethod']} {event['resource']}"
    try:
        fault.apply(SERVICE, fault.current(ssm, os.environ["FAULT_PARAM"]))
        if route == "POST /orders":
            res = create(json.loads(event['body']), ddb)
        elif route == "GET /orders/{id}":
            res = get(event['pathParameters']['id'], ddb)
        elif route == "POST /orders/{id}/cancel":
            res = cancel(event['pathParameters']['id'], ddb)
        else:
            res = respond(404, {"error": f"no route {route}"})
    except calls.DownstreamTimeout as exc:
        res = respond(504, {"error": f"dependency timed out: {exc}"})
    except calls.DownstreamError as exc:
        res = respond(502, {"error": f"dependency failed: {exc}"})
    except fault.InjectedFailure as exc:
        res = respond(500, {"error": str(exc)})
    level = "ERROR" if res["statusCode"] >= 500 else "INFO"
    obs.log(level, "request", route=route, status=res["statusCode"],
            ms=(time.perf_counter() - t0) * 1000)
    return res
\end{verbatim}

\textbf{Key design choices:}
\begin{itemize}
    \item Fault switch checked on every invocation (cached 5s in \texttt{fault.current})
    \item Downstream timeouts (504) distinguished from errors (502)
    \item Structured JSON logging with \texttt{obs.log} for CloudWatch
    \item X-Ray context automatically propagated by AWS SDK
\end{itemize}

\subsection{inventory Handler (Downstream Service)}

\begin{verbatim}
def handler(event, context):
    ddb, ssm = clients()
    fault.apply(SERVICE, fault.current(ssm, os.environ["FAULT_PARAM"]))
    action = event.get("action")
    if action == "reserve":
        return reserve(event["sku"], event["qty"], ddb)
    return {"error": "unknown action"}

def reserve(sku, qty, ddb):
    item = ddb.get_item(TableName=os.environ["TABLE_NAME"],
                        Key={"pk": {"S": f"SKU#{sku}"}}).get("Item")
    if not item or int(item["stock"]["N"]) < qty:
        return {"reserved": False}
    return {"reserved": True, "price_cents": int(item["price_cents"]["N"])}
\end{verbatim}

\textbf{Fault application:} \texttt{fault.apply} is called before business logic. If \texttt{service == target}, the configured fault (delay, exception, etc.) is applied.

\section{Fault Injection Implementation}

\subsection{Schedule Generation}

\texttt{injector/schedule.py} generates the injection timeline:

\begin{verbatim}
def generate_schedule(fault_types, targets, load_levels,
                      injections_per_type, seed=42):
    rng = np.random.default_rng(seed)
    schedule = []
    for fault_type in fault_types:
        for load_level in load_levels:
            for i in range(injections_per_type):
                target = targets[i % len(targets)]  # Round-robin
                schedule.append({
                    "id": len(schedule),
                    "type": fault_type,
                    "target": target,
                    "load": load_level,
                    "duration_s": 60,
                    "recovery_s": 300
                })
    rng.shuffle(schedule)  # Randomise order
    return schedule
\end{verbatim}

\textbf{Output:} A JSON list of 480 injections, seeded for reproducibility.

\subsection{Injector Execution}

\texttt{injector/injector.py} writes faults to SSM and applies throttling:

\begin{verbatim}
def inject(ssm, lam, stack, fault):
    if fault["type"] == "throttling":
        # Set reserved concurrency to 0 = all invocations throttled
        lam.put_function_concurrency(FunctionName=f"{stack}-{fault['target']}",
                                      ReservedConcurrentExecutions=0)
    # Write fault JSON to SSM
    ssm.put_parameter(Name=f"/{stack}/fault",
                      Value=json.dumps({
                          "id": fault["id"],
                          "type": fault["type"],
                          "target": fault["target"],
                          "until": time.time() + fault["duration_s"],
                          "latency_ms": fault.get("latency_ms", 500),
                          "hold_ms": fault.get("hold_ms", 1500)
                      }),
                      Overwrite=True)
    # Log ground truth
    with open(f"{run_dir}/ground_truth.jsonl", "a") as f:
        f.write(json.dumps({...}) + "\n")
\end{verbatim}

\textbf{Throttling special case:} Injecting \texttt{dependency\_failure} via SSM works for most faults, but throttling requires AWS API call to modify reserved concurrency.

\section{Detector Implementation}

\subsection{CloudWatch Rule-Based Detector}

\texttt{detector/rules.py} implements the 3-sigma + P99 detection:

\begin{verbatim}
def detect(metrics, thresholds, window_minutes=30):
    alarms = []
    for service, service_metrics in metrics.groupby("service"):
        for metric_name in ["Errors", "Duration", "Throttles"]:
            recent = service_metrics.tail(window_minutes)
            mu, sigma = recent[metric_name].mean(), recent[metric_name].std()
            current = recent[metric_name].iloc[-1]
            threshold = thresholds[service][metric_name]
            if current > mu + 3 * sigma or current > threshold:
                alarms.append({
                    "time": recent["timestamp"].iloc[-1],
                    "service": service,
                    "metric": metric_name,
                    "value": current,
                    "threshold": threshold
                })
    return alarms
\end{verbatim}

\textbf{Threshold enforcement:} The detector verifies that \texttt{thresholds} hash matches the frozen calibration config; refuses to run otherwise.

\subsection{X-Ray Dependency Ranker}

\texttt{detector/ranker.py} implements the trace-based ranking:

\begin{verbatim}
def rank_services(traces, p99_thresholds):
    failed_traces = [t for t in traces if is_failed(t, p99_thresholds)]
    service_scores = {}
    for service in get_all_services(traces):
        fail_share = sum(1 for t in failed_traces if service in t) / len(failed_traces)
        depths = [get_depth(t, service) for t in traces if service in t]
        depth_score = 1 / np.mean(depths) if depths else 0
        service_scores[service] = 0.5 * fail_share + 0.5 * depth_score
    return sorted(service_scores.items(), key=lambda x: -x[1])

def is_failed(trace, p99_thresholds):
    for span in trace["spans"]:
        if span.get("error"):
            return True
        if span["self_time_ms"] > p99_thresholds.get(span["name"], float("inf")):
            return True
    return False
\end{verbatim}

\textbf{Span P99 calibration:} During the 24-hour calibration, we compute P99 self-time for each service's spans. This allows detection of \texttt{elevated\_latency} faults even when no error is raised.

\section{Evaluation Pipeline Implementation}

\subsection{Rig Evaluation}

\texttt{eval/rig.py} orchestrates end-to-end evaluation:

\begin{verbatim}
def run_rig_evaluation(source, run_dir, out_dir, figures_dir):
    # 1. Load ground truth
    ground_truth = load_ground_truth(f"{run_dir}/ground_truth.jsonl")
    
    # 2. Collect telemetry
    if source == "sim":
        metrics, traces = sim.generate_telemetry(ground_truth)
    else:
        metrics = collectors.fetch_cloudwatch_metrics(run_dir)
        traces = collectors.fetch_xray_traces(run_dir)
    
    # 3. Calibrate thresholds (or load frozen)
    thresholds = load_frozen_thresholds("configs/experiment.yaml")
    
    # 4. Run detection
    detections = rules.detect(metrics, thresholds)
    
    # 5. Run localisation
    rankings = []
    for detection in detections:
        traces_window = filter_traces(traces, detection["time"])
        rank = ranker.rank_services(traces_window, thresholds["span_p99"])
        rankings.append(rank)
    
    # 6. Compute metrics
    results = metrics_module.evaluate_detection(ground_truth, detections)
    results.update(metrics_module.evaluate_localisation(ground_truth, rankings))
    
    # 7. Visualise
    plot_detection_localisation(results, f"{figures_dir}/rig.png")
    
    # 8. Write summary
    write_summary(results, f"{out_dir}/summary.md")
    return results
\end{verbatim}

\subsection{RCAEval Comparison}

\texttt{baseline\_runner/rule\_arm.py} adapts the rule-based method to RCAEval schema:

\begin{verbatim}
def evaluate_on_rcaeval(case):
    # RCAEval provides 600s before + 600s after injection
    metrics = case["metrics"]  # 10-second buckets
    traces = case["traces"]
    
    # Calibrate on first 300s (control period)
    calibration_metrics = metrics[:30]  # 30 buckets × 10s
    thresholds = compute_thresholds(calibration_metrics)
    
    # Detect on 300-600s (injection period)
    injection_metrics = metrics[30:60]
    detections = detect(injection_metrics, thresholds)
    
    # Localise using all 600s of traces
    if detections:
        rank = ranker.rank_services(traces, thresholds["span_p99"])
        return {"detected": True, "rank": rank}
    return {"detected": False, "rank": []}
\end{verbatim}

\textbf{Adaptation challenges:} RCAEval metrics are in 10-second buckets (vs 1-minute CloudWatch aggregation). We scale threshold floors proportionally to avoid spurious alarms on quantisation noise.

\section{Testing Strategy}

76 tests organised into:

\begin{itemize}
    \item \textbf{Handler tests (\texttt{test\_app.py}):} Local execution (FAULTLAB\_LOCAL=1) with moto-mocked DynamoDB/SSM
    \item \textbf{Fault injection tests (\texttt{test\_injector.py}):} Schedule generation, SSM writes, ground truth logging
    \item \textbf{Detector tests (\texttt{test\_rules.py}):} Threshold calibration, 3-sigma detection, hash verification
    \item \textbf{Ranker tests (\texttt{test\_ranker.py}):} Fail-share computation, depth scoring, span P99 checks
    \item \textbf{Metrics tests (\texttt{test\_metrics.py}):} P/R/F1 calculation, Top-k accuracy, delay quantiles
    \item \textbf{Integration tests (\texttt{test\_rig.py}):} End-to-end pipeline on simulated telemetry
    \item \textbf{SAM template tests (\texttt{test\_template.py}):} cfn-lint validation, resource naming
\end{itemize}

\textbf{Coverage:} pytest --cov reports 89\% line coverage (uncovered lines are AWS API error paths not exercisable in mocked environment).

\section{Continuous Integration}

Makefile targets automate testing and deployment:

\begin{verbatim}
make setup         # Create venv, install dependencies
make test          # Run pytest
make lint          # Ruff linting
make cfn-lint      # Validate SAM template
make check-bib     # Verify DOI resolution (offline)
make sim           # Run simulated rig evaluation
make build         # SAM build
make deploy        # SAM deploy (requires AWS credentials)
make calibration   # 24-hour threshold calibration
make steady        # Steady load campaign (2 RPS)
make peak          # Peak load campaign (10 RPS)
make collect       # Fetch CloudWatch/X-Ray telemetry
make rig           # Evaluate detection/localisation
\end{verbatim}

\section{Summary}

The implementation delivers a production-ready AWS serverless application with:
\begin{itemize}
    \item Declarative Infrastructure as Code (AWS SAM)
    \item Four Lambda functions with realistic order-processing logic
    \item Controlled fault injection via SSM Parameter Store
    \item Rule-based detection with cryptographically frozen thresholds
    \item X-Ray trace-based localisation with calibrated span thresholds
    \item Complete evaluation pipeline (rig + RCAEval + overhead)
    \item 76 passing tests with 89\% coverage
    \item Reproducible via Makefile from a clean environment
\end{itemize}

The next chapter presents the evaluation results: detection/localisation accuracy, baseline comparison, overhead measurements, and hypothesis test outcomes.
